% section: 等比数列の総和とその導出

\section{等比数列の総和とその導出}\label{節：等比数列の総和とその導出}
  \noindent\textbf{【本編で使う場面】}\par
  等比数列の和を,倍率をそろえて引く操作から導く.\sectionref*{節：階差型}と\sectionref*{節：特性方程式型}の階差による別解で使う.

  初項が$a$,公比が$r$である等比数列
  \[
    a,\,ar,\,ar^2,\,\ldots,\,ar^{n-1}
  \]
  の初項から第$n$項までの総和を考える.
  その総和を$S_n$とすると,
  \[
    S_n
    =\sum_{k=1}^{n} ar^{k-1}
    =a+ar+ar^2+\cdots+ar^{n-1}
  \]
  と表せる.

  \subsection{総和の公式}\label{節：等比数列の総和とその導出：総和の公式}
    公比$r$が$1$でないとき,
    \[
      S_n=\frac{a(1-r^n)}{1-r}
    \]
    である.
    分母と分子の符号を同時に変えれば,
    \[
      S_n=\frac{a(r^n-1)}{r-1}
    \]
    とも書ける.
    公比が$1$のときは,すべての項が$a$であるから,
    \[
      S_n=na
    \]
    となる.

  \subsection{公式の導出}\label{節：等比数列の総和とその導出：公式の導出}
    $r\ne1$として,総和の式の両辺に$r$を掛ける.
    \begin{align*}
      S_n
        &=a+ar+ar^2+\cdots+ar^{n-1},\\
      rS_n
        &=ar+ar^2+ar^3+\cdots+ar^n.
    \end{align*}
    2つの式を項ごとに引くと,右辺では共通する項がすべて消える.
\BookMathLayout{\[
    \begin{array}{rrrl}
        &S_{n}={}&a+&\hspace{-0.75em}ar+ar^2+\cdots+ar^{n-1}\\
      -\,) & rS_n={}&{}&\hspace{-0.75em}ar+ar^2+ar^3+\cdots+ar^{n}\\ \hline
        & (1-r)S_n={}&{}&\hspace{-0.75em}a-ar^n
    \end{array}
    \]}{
\[
\begin{array}{rl}
S_n={}&a+ar+ar^2+\cdots\\
     &{}+ar^{n-1}\\
-\,)\ rS_n={}&ar+ar^2+ar^3+\cdots\\
     &{}+ar^n\\ \hline
(1-r)S_n={}&a-ar^n
\end{array}
\]
}
    よって,
    \[
      (1-r)S_n=a(1-r^n)
    \]
    である.
    $r\ne1$より両辺を$1-r$で割ると,
    \[
      S_n=\frac{a(1-r^n)}{1-r}
    \]
    を得る.

    この導出では,総和を公比$r$倍してから元の総和との差を取ることで,途中の項を消している.
    数列の和を求めるときには,このように総和を定数倍して引き,共通する項を消す考え方がよく用いられる.

